\documentclass[11pt, a4paper]{article}

%% ── Packages ──────────────────────────────────────────────────────────────
\usepackage[T1]{fontenc}
\usepackage[utf8]{inputenc}
\usepackage{lmodern}
\usepackage{microtype}
\usepackage[margin=2.5cm]{geometry}
\usepackage{amsmath, amssymb, amsthm}
\usepackage{booktabs, longtable, array, multirow}
\usepackage{xcolor}
\usepackage{hyperref}
\usepackage{listings}
\usepackage{graphicx}
\usepackage{float}
\usepackage{enumitem}
\usepackage{parskip}
\usepackage{titlesec}
\usepackage{fancyhdr}
\usepackage{caption}

%% ── Colours ───────────────────────────────────────────────────────────────
\definecolor{cblue}{RGB}{25,  115, 217}
\definecolor{cgray}{RGB}{90,  90,  90}
\definecolor{cred} {RGB}{200, 50,  50}
\definecolor{cgreen}{RGB}{30, 140, 60}
\definecolor{clightgray}{RGB}{245,245,245}
\definecolor{codegray}{RGB}{100,100,100}

%% ── Listings (MATLAB) ─────────────────────────────────────────────────────
\lstdefinestyle{matlab}{
  backgroundcolor=\color{clightgray},
  basicstyle=\ttfamily\small,
  keywordstyle=\color{cblue}\bfseries,
  commentstyle=\color{codegray}\itshape,
  stringstyle=\color{cgreen},
  breaklines=true,
  captionpos=b,
  frame=single,
  framesep=4pt,
  numbers=none,
  tabsize=4,
  morekeywords={function,end,for,if,else,elseif,while,try,catch,
                arguments,timetable,table,datetime,true,false,NaN,Inf},
}
\lstset{style=matlab}

%% ── Headers / footers ─────────────────────────────────────────────────────
\pagestyle{fancy}
\fancyhf{}
\lhead{\small\textcolor{cgray}{Two-Step Equity Underweight Strategy}}
\rhead{\small\textcolor{cgray}{\thepage}}
\renewcommand{\headrulewidth}{0.4pt}

%% ── Hyperref ──────────────────────────────────────────────────────────────
\hypersetup{
  colorlinks=true,
  linkcolor=cblue,
  urlcolor=cblue,
  citecolor=cblue,
  pdftitle={Two-Step Equity Underweight Strategy: Technical Documentation},
  pdfauthor={CIO Quantitative Research},
}

%% ── Title ─────────────────────────────────────────────────────────────────
\title{%
  \vspace{-1cm}
  {\Large\bfseries Two-Step Equity Underweight Strategy}\\[0.4em]
  {\large Technical Model Documentation}\\[0.8em]
  {\normalsize\textcolor{cgray}{Version 1.0 \textbar{} CIO Quantitative Research}}
}
\author{}
\date{\today}

%% ══════════════════════════════════════════════════════════════════════════
\begin{document}
\maketitle
\thispagestyle{empty}
\vspace{-1em}
\noindent\rule{\linewidth}{0.8pt}

%% ── Abstract ──────────────────────────────────────────────────────────────
\begin{abstract}
\noindent
This document provides a complete technical description of the Two-Step Equity
Underweight Strategy. The model monitors four macro and market signals
(payrolls momentum, yield-curve slope, HY credit spreads, and equity price
trend) and dynamically reduces equity allocation from a neutral \textbf{50\%}
to a mild-underweight \textbf{35\%} or severe-underweight \textbf{20\%}
when the composite signal falls below calibrated thresholds.

Over the full backtested period (December 1996 -- November 2025, 29 years),
the strategy delivered an annualised return of \textbf{8.14\%} at
\textbf{5.87\%} annualised volatility, versus 7.97\% at 15.38\% for a fully
invested equity position and 6.04\% at 6.79\% for a matched-exposure passive
benchmark. The maximum drawdown was reduced from $-52.6\%$ (100\% equities)
to $-10.7\%$, with an annualised Sharpe ratio of \textbf{1.39}.
\end{abstract}

\noindent\rule{\linewidth}{0.4pt}
\tableofcontents
\clearpage

%% ══════════════════════════════════════════════════════════════════════════
\section{Strategy Overview}
\label{sec:overview}
%% ══════════════════════════════════════════════════════════════════════════

The strategy is a rules-based, signal-driven overlay applied to a static
mixed portfolio. The neutral allocation ($w_N = 50\%$ equity, $50\%$
investment-grade bonds) is maintained unless the composite market signal
crosses a calibrated threshold, at which point equity is reduced and the
freed allocation is placed in USD cash.

\subsection{Allocation States}

\begin{table}[H]
\centering
\caption{Allocation states and portfolio composition}
\label{tab:alloc-states}
\begin{tabular}{lcccc}
\toprule
\textbf{State} & \textbf{Trigger Condition} & \textbf{Equity} & \textbf{Bonds} & \textbf{Cash} \\
\midrule
Neutral   & $M_t > \theta_m$                           & 50\% & 50\% & 0\%  \\
Mild UW   & $\theta_s < M_t \leq \theta_m$             & 35\% & 50\% & 15\% \\
Severe UW & $M_t \leq \theta_s$                        & 20\% & 50\% & 30\% \\
\bottomrule
\end{tabular}
\end{table}

where $M_t$ denotes the \emph{Market Composite} signal (Section~\ref{sec:mktcomp})
and $(\theta_m, \theta_s)$ are the mild and severe signal thresholds
(Section~\ref{sec:threshold}). The bond allocation is held constant at 50\%
across all states. The de-risked equity portion earns the prevailing Federal
Funds target rate.

\subsection{Pipeline Architecture}

The MATLAB implementation follows a strict
Model--View--Controller (MVC) architecture using MATLAB packages:

\begin{description}[leftmargin=3em, labelindent=1em]
  \item[\texttt{+Model/}] Signal generation, data access, backtesting logic.
  \item[\texttt{+View/}]  GUI application (\texttt{UnderweightApp}).
  \item[\texttt{+Controller/}] Orchestration (\texttt{UnderweightController}).
  \item[\texttt{Underweight/}] Strategy scripts and helper functions.
\end{description}

The core execution pipeline is:
\[
  \text{buildSignals} \;\longrightarrow\; \text{labelEquityRegimes}
  \;\longrightarrow\; \text{thresholdApproach}
  \;\longrightarrow\; \text{evalPortfolio}
\]

%% ══════════════════════════════════════════════════════════════════════════
\section{Data Sources}
\label{sec:data}
%% ══════════════════════════════════════════════════════════════════════════

\begin{table}[H]
\centering
\caption{Data sources used in the strategy}
\label{tab:data-sources}
\begin{tabular}{llll}
\toprule
\textbf{Series} & \textbf{Provider} & \textbf{Frequency} & \textbf{Available from} \\
\midrule
S\&P 500 Index (\texttt{\^{}SPX}) & Stooq & Daily OHLCV   & January 1928 \\
PAYEMS (Nonfarm Payrolls)    & FRED  & Monthly (SA)  & January 1939 \\
T10Y2Y (10Y$-$2Y Spread)     & FRED  & Daily (\%)    & June 1976 \\
BAMLH0A0HYM2 (HY OAS)        & FRED  & Daily (\%)    & December 1996 \\
FEDFUNDS (Effective FFR)     & FRED  & Monthly (\%)  & July 1954 \\
\bottomrule
\end{tabular}
\end{table}

All data is fetched via the \texttt{Model.DataManager} class, which
maintains an on-disk SQLite registry and \texttt{.mat} cache. Incremental
updates are performed automatically: a fresh request with no \texttt{startDate}
only downloads observations newer than the last cached date.

The equity price priority chain is:
Yahoo Finance (rate-limited) $\rightarrow$ Stooq SPY (data from 2005)
$\rightarrow$ Stooq \texttt{\^{}SPX} (data from 1994). The fallback to
\texttt{\^{}SPX} is triggered if the fetched history is shorter than 25 years
($252 \times 25$ trading days).

The effective backtest start date is constrained by BAMLH0A0HYM2
(December 1996), though data is fetched from January 1994 to provide a
60-month warm-up for rolling z-score calculations.

%% ══════════════════════════════════════════════════════════════════════════
\section{Signal Construction}
\label{sec:signals}
%% ══════════════════════════════════════════════════════════════════════════

Each signal maps to the interval $(-1, +1)$ with the convention:
\[
  s_t > 0 \;\Leftrightarrow\; \textit{risk-on}, \qquad
  s_t < 0 \;\Leftrightarrow\; \textit{risk-off}.
\]

Monthly signal values are obtained by taking the last non-NaN observation
within each calendar month (end-of-month resampling).

%% ── 3.1 Payrolls Momentum ────────────────────────────────────────────────
\subsection{Payrolls Momentum (\texttt{PayrollsMom})}
\label{sec:payrolls}

\textbf{Source:} FRED PAYEMS (Total Nonfarm Payrolls, thousands, monthly SA).

The signal is constructed via a five-step pipeline that provides a noise-robust,
adaptive momentum indicator for the US labour market.

\paragraph{Step 1 — Raw month-over-month change}
\[
  \Delta P_t = P_t - P_{t-1}
\]

\paragraph{Step 2 — Trailing median prefilter (noise gate)}
\[
  \widetilde{\Delta P}_t = \operatorname{median}\!\bigl(\Delta P_{t-W+1},\,\ldots,\,\Delta P_t\bigr),
  \quad W = 3
\]
A single anomalous print (weather event, strike, seasonal revision) moves the
trailing median of three to the middle value, effectively replacing the outlier
with a neighbour. Two or more consecutive prints in the same direction pass
through unchanged, preserving genuine labour-market inflections with at most
one month of additional lag.

\paragraph{Step 3 — Exponential moving average (numerator path)}
\[
  M_t = \alpha\,\widetilde{\Delta P}_t + (1-\alpha)\,M_{t-1},
  \quad \alpha = 1 - e^{-\ln 2\,/\,h},\quad h = 3\ \text{months}
\]

\paragraph{Step 4 — Expanding-window winsorisation (denominator path)}
\[
  Q_t = \operatorname{quantile}\bigl(|\Delta P_1|,\ldots,|\Delta P_t|;\,p=95\bigr),
  \qquad
  \widehat{\Delta P}_t = \operatorname{sign}(\Delta P_t)\cdot\min\!\bigl(|\Delta P_t|, Q_t\bigr)
\]
This prevents catastrophic single-month shocks (e.g.\ $-$20{,}700~K in April
2020 due to COVID-19) from inflating the rolling standard deviation and
suppressing the signal for subsequent years.

\paragraph{Step 5 — Adaptive normalisation and tanh compression}
\[
  \sigma_t = \operatorname{std}\!\bigl(\widehat{\Delta P}_{t-N+1},\ldots,\widehat{\Delta P}_t\bigr),
  \quad N = 60
\]
\[
  \boxed{s_t^{\text{Pay}} = \tanh\!\left(\frac{M_t}{\sigma_t}\right) \in (-1,+1)}
\]

\textit{Calibration:} Under default parameters, a sustained expansion
($+200$~K/month) gives $s_t \approx +0.80$; a mild recession ($-100$~K/month)
gives $s_t \approx -0.59$; an isolated shock is filtered to $\approx 0$.

%% ── 3.2 Yield Curve Slope ────────────────────────────────────────────────
\subsection{Yield Curve Slope (\texttt{YieldCurve})}
\label{sec:yield}

\textbf{Source:} FRED T10Y2Y (10-Year minus 2-Year Treasury Constant Maturity
Rate, daily, percentage points).

\paragraph{Step 1 — Raw spread}
$S_t$ = T10Y2Y level (percentage points). Positive = normal (risk-on);
negative = inverted (risk-off).

\paragraph{Step 2 — Causal rolling z-score}
\[
  z_t = \frac{S_t - \mu_t}{\sigma_t},
  \quad
  \mu_t = \operatorname{mean}(S_{t-W_d+1},\ldots,S_t),
  \quad
  \sigma_t = \operatorname{std}(S_{t-W_d+1},\ldots,S_t)
\]
where $W_d = W \times 21$ trading days, $W = 60$ months. A steep curve above
its historical average yields $z > 0$; an inverted curve below its average
yields $z < 0$.

\paragraph{Step 3 — EMA smoothing}
\[
  m_t = \alpha\, z_t + (1-\alpha)\, m_{t-1},
  \quad \alpha = 1 - e^{-\ln 2\,/\,h},\quad h = 3\ \text{months}
\]

\paragraph{Step 4 — Tanh compression}
\[
  \boxed{s_t^{\text{YC}} = \tanh(m_t) \in (-1,+1)}
\]

%% ── 3.3 Credit Spread Signal ─────────────────────────────────────────────
\subsection{Credit Spread Signal (\texttt{CreditSpread})}
\label{sec:credit}

\textbf{Source:} FRED BAMLH0A0HYM2 (ICE BofA US High Yield Index
Option-Adjusted Spread, daily, \%).

The signal is \emph{inverted}: wide spreads indicate credit stress and
produce a negative (risk-off) output.

\paragraph{Step 1 — Raw OAS}
$C_t$ = daily HY OAS level (\%).

\paragraph{Step 2 — Expanding-window winsorisation (denominator path)}
\[
  Q_t = \operatorname{quantile}(C_1,\ldots,C_t;\;p=97),
  \quad \hat{C}_t = \min(C_t, Q_t)
\]
This caps the denominator series against extreme spike events (GFC 2008:
$\approx 20\%$; COVID 2020: $\approx 11\%$) that would otherwise inflate
$\sigma_t$ and flatten the signal for years.

\paragraph{Step 3 — Causal rolling z-score}
\[
  z_t = \frac{C_t - \mu_t(\hat{C})}{\sigma_t(\hat{C})},
  \quad W = 60\ \text{months}
\]

\paragraph{Step 4 — EMA smoothing}
\[
  m_t = \alpha\, z_t + (1-\alpha)\, m_{t-1},
  \quad \alpha = 1 - e^{-\ln 2\,/\,h},\quad h = 1\ \text{month}
\]

\paragraph{Step 5 — Tanh compression, inverted}
\[
  \boxed{s_t^{\text{CS}} = \tanh(-m_t) \in (-1,+1)}
\]

The short EMA half-life ($h = 1$) ensures the signal reacts immediately to
credit market blowouts, which are typically fast-moving and highly predictive
of equity drawdowns.

%% ── 3.4 Equity Trend Signal ──────────────────────────────────────────────
\subsection{Equity Trend Signal (\texttt{EquityTrend})}
\label{sec:trend}

\textbf{Source:} S\&P 500 daily close price $P_t$ (Stooq \texttt{\^{}SPX}).

\paragraph{Step 1 — Trailing moving average}
\[
  \operatorname{MA}_t = \frac{1}{N}\sum_{i=0}^{N-1} P_{t-i},
  \quad N = 200\ \text{days}
\]
NaN is assigned for the first $N-1$ observations (warm-up period).

\paragraph{Step 2 — Tanh-compressed deviation}
\[
  \boxed{s_t^{\text{ET}} = \tanh\!\left(k\,\Bigl(\frac{P_t}{\operatorname{MA}_t} - 1\Bigr)\right) \in (-1,+1),
  \quad k = 10}
\]

\textit{Calibration:} 10\% above the 200-day MA gives
$\tanh(10 \times 0.10) = \tanh(1) \approx 0.76$; 10\% below gives $\approx -0.76$.
The scale $k = 10$ provides a graded signal rather than a hard step function.

%% ── 3.5 Equal-Weighted Composite ─────────────────────────────────────────
\subsection{Equal-Weighted Composite (\texttt{Composite})}
\label{sec:composite}

\[
  \boxed{C_t = \frac{1}{n_t}\sum_{i \in \mathcal{S}} s_{i,t},
  \quad n_t = \bigl|\{i \in \mathcal{S} : s_{i,t} \neq \text{NaN}\}\bigr|}
\]

where $\mathcal{S} = \{\text{PayrollsMom},\,\text{YieldCurve},\,\text{CreditSpread},\,\text{EquityTrend}\}$.
The NaN-tolerant averaging ensures the composite is defined even during the
MA warm-up period and any data gaps.

%% ── 3.6 Market Composite ──────────────────────────────────────────────────
\subsection{Market Composite (\texttt{MarketComposite})}
\label{sec:mktcomp}

A weighted composite that over-weights the two market-price signals
(CreditSpread and EquityTrend) relative to the macro signal (YieldCurve):

\[
  \boxed{M_t = \frac{w_{\text{CS}}\,s_t^{\text{CS}}
              + w_{\text{ET}}\,s_t^{\text{ET}}
              + w_{\text{YC}}\,s_t^{\text{YC}}}
             {\sum_{i} w_i \cdot \mathbf{1}[s_{i,t} \neq \text{NaN}]}}
\]

\begin{table}[H]
\centering
\caption{MarketComposite signal weights}
\begin{tabular}{lcc}
\toprule
\textbf{Component} & \textbf{Series} & \textbf{Weight} \\
\midrule
Credit Spread & BAMLH0A0HYM2 & 30\% \\
Equity Trend  & S\&P 500 vs 200d MA & 50\% \\
Yield Curve   & T10Y2Y & 20\% \\
\midrule
\textbf{Total} & & \textbf{100\%} \\
\bottomrule
\end{tabular}
\end{table}

The denominator re-normalises to the sum of \emph{available} weights,
so the composite remains well-defined if any component is NaN.
PayrollsMom is excluded from MarketComposite because its monthly frequency
and lagged publication make it redundant relative to the higher-frequency
market-price signals.

\subsection{Monthly Resampling}
\label{sec:resample}

All daily signals are resampled to a common monthly end-of-month grid:
\[
  s_{\text{month}} = \text{last non-NaN observation in calendar month}
\]
The grid is anchored to the latest start month across all four series
(December 1996, constrained by BAMLH0A0HYM2) and extends to the most
recent complete month. A year$\times$100\,+\,month integer key is used
for timezone-agnostic month matching.

%% ══════════════════════════════════════════════════════════════════════════
\section{Regime Classification}
\label{sec:labels}
%% ══════════════════════════════════════════════════════════════════════════

Forward-looking regime labels are assigned at each month $t$ based on the
subsequent three months of equity returns and realised volatility. These
labels are used \emph{only} for threshold calibration and decision-tree
training; they are never used in live signal generation.

Let $r_{t+k}^{\text{eq}}$ denote the equity log-return for month $t+k$.
Define the 3-month forward cumulative return and annualised realised
volatility:
\[
  R_t^{3m} = \prod_{k=1}^{3}\bigl(1 + r_{t+k}^{\text{eq}}\bigr) - 1,
  \qquad
  \sigma_t^{3m} = \operatorname{std}\bigl(r_{t+1}^{\text{eq}},\,r_{t+2}^{\text{eq}},\,r_{t+3}^{\text{eq}}\bigr)\times\sqrt{12}
\]

The label assigned to month $t$ is:
\[
  \ell_t =
  \begin{cases}
    2 & \text{(Severe UW)} \quad \text{if } R_t^{3m} < 0 \;\text{ and }\; \sigma_t^{3m} > 30\% \\
    1 & \text{(Mild UW)}   \quad \text{if } R_t^{3m} < 0 \;\text{ and }\; \sigma_t^{3m} \leq 30\% \\
    0 & \text{(Neutral)}   \quad \text{otherwise} \\
    \text{NaN} & \text{if fewer than 3 forward months available}
  \end{cases}
\]

The intuition is that high-volatility negative-return periods are the
most destructive for long-term portfolio compounding and justify the
most aggressive underweighting (Severe). Low-volatility negative-return
periods (Mild) justify a moderate reduction.

\subsection{Sample Distribution}

Over the backtested period (December 1996 -- November 2025, 348 labelled months):

\begin{table}[H]
\centering
\caption{Regime label distribution}
\label{tab:labels}
\begin{tabular}{lrrl}
\toprule
\textbf{Regime} & \textbf{Months} & \textbf{\%} & \textbf{Allocation} \\
\midrule
Neutral         & 242 & 69.5\% & 50\% equity \\
Mild UW         &  87 & 25.0\% & 35\% equity \\
Severe UW       &  19 &  5.5\% & 20\% equity \\
\midrule
\textbf{Total}  & 348 & 100\%  & \\
\bottomrule
\end{tabular}
\end{table}

The severe class represents only 5.5\% of observations, corresponding
primarily to the GFC (2008--2009), COVID crash (2020), and the 2022 rate-hike
bear market.

%% ══════════════════════════════════════════════════════════════════════════
\section{Signal Threshold Optimisation}
\label{sec:threshold}
%% ══════════════════════════════════════════════════════════════════════════

\subsection{Methodology}

For each signal or composite $s_t$, the threshold pair
$(\hat{\theta}_m, \hat{\theta}_s)$ is found by exhaustive grid search:

\[
  (\hat{\theta}_m, \hat{\theta}_s) = \arg\max_{\theta_s \leq \theta_m}
  \operatorname{BalAcc}\!\bigl(a(s, \theta_m, \theta_s),\, \ell\bigr)
\]

where the predicted allocation $a_t$ maps to a label via:
\[
  \hat{\ell}_t(\theta_m, \theta_s) =
  \begin{cases}
    2 & s_t \leq \theta_s \\
    1 & \theta_s < s_t \leq \theta_m \\
    0 & s_t > \theta_m
  \end{cases}
\]

and balanced accuracy is the mean per-class recall:
\[
  \operatorname{BalAcc} = \frac{1}{3}\left(
    \frac{\mathrm{TP}_0}{|\ell=0|} +
    \frac{\mathrm{TP}_1}{|\ell=1|} +
    \frac{\mathrm{TP}_2}{|\ell=2|}
  \right)
\]

Balanced accuracy is used (rather than overall accuracy) to account for
the severe class imbalance: a na\"{i}ve classifier that always predicts
Neutral achieves 69.5\% overall accuracy but BalAcc of only 0.333.

The threshold grids used are:
\[
  \Theta_m \in \{-0.9, -0.8, \ldots, +0.9\},\quad
  \Theta_s \in \{-0.9, -0.8, \ldots, +0.9\}
\]
with the constraint $\theta_s \leq \theta_m$, giving approximately 171 valid
$(\theta_m, \theta_s)$ pairs per signal.

\subsection{Results}

\begin{table}[H]
\centering
\caption{Optimised thresholds and balanced accuracy by signal variant}
\label{tab:thresholds}
\begin{tabular}{lrrrr}
\toprule
\textbf{Signal}       & \textbf{Mild thr} $\hat{\theta}_m$ & \textbf{Severe thr} $\hat{\theta}_s$ & \textbf{BalAcc} & \textbf{Avg Alloc} \\
\midrule
PayrollsMom      & $-0.10$ & $-0.30$ & 0.498 & 45.9\% \\
YieldCurve       & $+0.10$ & $-0.80$ & 0.261 & 35.4\% \\
CreditSpread     & $-0.50$ & $-0.60$ & 0.469 & 41.2\% \\
EquityTrend      & $+0.10$ & $-0.40$ & 0.510 & 43.1\% \\
Composite        & $0.00$  & $-0.20$ & 0.499 & 43.9\% \\
MacroComposite   & $0.00$  & $-0.30$ & 0.489 & 45.1\% \\
\textbf{MarketComposite} & $\mathbf{-0.20}$ & $\mathbf{-0.40}$ & \textbf{0.492} & \textbf{44.1\%} \\
\bottomrule
\end{tabular}
\end{table}

The YieldCurve signal achieves the lowest balanced accuracy (0.261) because
the yield curve is a slow-moving structural indicator; it leads equity
drawdowns by 12--18 months on average, making near-term classification
difficult. The EquityTrend signal achieves the highest individual BalAcc
(0.510), reflecting the reflexive relationship between price momentum and
regime.

The \textbf{MarketComposite} is selected as the primary signal for live
allocation because it:
\begin{enumerate}[leftmargin=2em]
  \item Combines both market-price signals (which lead on a 1--3 month horizon)
        with the macro yield signal (which provides leading information on
        recessions);
  \item Achieves a competitive BalAcc of 0.492 at an average equity exposure
        of 44.1\%;
  \item Is robust to parameter perturbation given that its components react
        on different timescales.
\end{enumerate}

%% ══════════════════════════════════════════════════════════════════════════
\section{Decision Tree Approach}
\label{sec:tree}
%% ══════════════════════════════════════════════════════════════════════════

As a secondary allocation method, a shallow classification tree
is trained on all four signal features to predict the regime label.
This approach allows non-linear interaction between signals without
requiring explicit composite construction.

\subsection{Model Specification}

\begin{description}[leftmargin=3em]
  \item[Algorithm] CART (\texttt{fitctree}, Statistics and Machine Learning
        Toolbox).
  \item[Features] $\mathbf{x}_t = [s_t^{\text{Pay}},\, s_t^{\text{YC}},\,
        s_t^{\text{CS}},\, s_t^{\text{ET}}]^\top$.
  \item[Target] Regime label $\ell_t \in \{0, 1, 2\}$.
  \item[Max splits] 4 (constrains tree to at most 5 leaves for interpretability).
  \item[Cost matrix] Asymmetric: misclassifying Severe as Neutral costs 5;
        misclassifying Mild as Neutral costs 2; symmetric otherwise. This
        penalises failure to de-risk during true stress regimes.
  \item[Validation] 5-fold stratified cross-validation.
\end{description}

\subsection{Trained Tree Rules}

The tree trained on the full December 1996 -- November 2025 sample yields
the following rules:

\begin{enumerate}[leftmargin=2em]
  \item If $s_t^{\text{Pay}} < -0.318$ and $s_t^{\text{ET}} < -0.396$
        $\Rightarrow$ \textbf{Severe UW} (20\% equity)
  \item If $s_t^{\text{Pay}} < -0.318$ and $s_t^{\text{ET}} \geq -0.396$
        $\Rightarrow$ \textbf{Neutral} (50\% equity)
  \item If $s_t^{\text{Pay}} \geq -0.318$ and $s_t^{\text{YC}} < -0.953$
        $\Rightarrow$ \textbf{Mild UW} (35\% equity)
  \item If $s_t^{\text{Pay}} \geq -0.318$ and $s_t^{\text{YC}} \geq -0.953$
        $\Rightarrow$ \textbf{Neutral} (50\% equity)
\end{enumerate}

\textit{Economic interpretation:} Severe underweighting requires simultaneous
deterioration in \emph{both} labour markets (PayrollsMom below $-0.318$)
\emph{and} equity price momentum (EquityTrend below $-0.396$). Mild
underweighting is triggered by a deeply inverted yield curve alone
(YieldCurve below $-0.953$), consistent with yield-curve inversion as an
intermediate-horizon recession signal.

\subsection{Performance}

\begin{table}[H]
\centering
\caption{Decision tree classification performance}
\begin{tabular}{lr}
\toprule
\textbf{Metric} & \textbf{Value} \\
\midrule
In-sample balanced accuracy    & 0.552 \\
5-fold CV balanced accuracy    & 0.391 \\
\midrule
\multicolumn{2}{l}{\textit{Confusion matrix (rows = actual, cols = predicted):}} \\
\midrule
Neutral correctly classified   & 232 / 242 \\
Mild correctly classified      &  10 / 87  \\
Severe correctly classified    &  11 / 19  \\
\bottomrule
\end{tabular}
\end{table}

The CV accuracy of 0.391 vs.\ in-sample 0.552 indicates moderate overfitting,
primarily from the under-represented Severe class (19 observations). The tree
is offered as a supplemental tool; the \emph{threshold approach} on
MarketComposite is the primary allocation engine.

%% ══════════════════════════════════════════════════════════════════════════
\section{Portfolio Construction and Evaluation}
\label{sec:portfolio}
%% ══════════════════════════════════════════════════════════════════════════

\subsection{Return Calculation}

At each month $t$ in the valid labelled subset, the strategy portfolio
return is:
\[
  r_t^{\text{port}} = a_t\, r_t^{\text{eq}}
                    + w_B\, r^{\text{bond}}
                    + (w_B - a_t)\, r_t^{\text{cash}}
\]
where:
\begin{itemize}[leftmargin=2em]
  \item $a_t \in \{0.20, 0.35, 0.50\}$ = signal-implied equity allocation;
  \item $w_B = 0.50$ = constant bond weight;
  \item $r^{\text{bond}} = (1 + 0.04)^{1/12} - 1 \approx 0.327\%$/month = assumed 4\%~p.a.;
  \item $r_t^{\text{cash}} = (1 + \text{FEDFUNDS}_t)^{1/12} - 1$ = monthly USD cash return
        from FRED FEDFUNDS (monthly average effective FFR).
\end{itemize}

Note that $w_B - a_t \geq 0$ for all valid states (since $a_t \leq w_B = 0.50$),
so the total portfolio weight always sums to 1.00:
\[
  a_t + w_B + (w_B - a_t) = a_t + 0.50 + 0.50 - a_t = 1.00\quad\checkmark
\]

Portfolio returns at NaN (warm-up) months are treated as zero.

\subsection{Benchmark Comparison}

Three strategies are compared:

\begin{description}[leftmargin=3em]
  \item[100\% Equities] $r_t = r_t^{\text{eq}}$. No risk management; full
        equity exposure throughout.
  \item[Matched Benchmark] $r_t = \bar{a}\, r_t^{\text{eq}} + (1-\bar{a})\,
        r^{\text{bond}}$, where $\bar{a} = 44.1\%$ is the time-average equity
        allocation of the Two-Step strategy. No timing; purely passive at the
        same long-run equity exposure.
  \item[Two-Step De-Risk] The signal-driven strategy as described above.
\end{description}

The Matched Benchmark is the economically appropriate comparator because
it eliminates the size effect: both strategies hold the same \emph{average}
equity position over time. Any excess return of the Two-Step strategy above
the Matched Benchmark is attributable purely to the timing ability of the
MarketComposite signal.

\subsection{Performance Metrics}

\begin{align}
  \hat{r}_{\text{ann}} &= \left(\prod_{t=1}^{T}(1+r_t)\right)^{12/T} - 1
  && \text{(annualised geometric return)} \\[4pt]
  \hat{\sigma}_{\text{ann}} &= \hat{\sigma}_m \sqrt{12}
  && \text{(annualised volatility)} \\[4pt]
  SR &= \hat{r}_{\text{ann}} / \hat{\sigma}_{\text{ann}}
  && \text{(annualised Sharpe ratio)} \\[4pt]
  \text{MDD} &= \min_{t}\frac{V_t - \max_{s \leq t} V_s}{\max_{s \leq t} V_s}
  && \text{(maximum peak-to-trough drawdown)}
\end{align}

%% ══════════════════════════════════════════════════════════════════════════
\section{Empirical Results}
\label{sec:results}
%% ══════════════════════════════════════════════════════════════════════════

\subsection{Sample and Allocation Statistics}

\begin{table}[H]
\centering
\caption{Allocation state dwell statistics (December 1996 -- November 2025)}
\label{tab:dwell}
\begin{tabular}{lrrrr}
\toprule
\textbf{State} & \textbf{Months} & \textbf{\%} & \textbf{Avg dwell} & \textbf{No. episodes} \\
\midrule
Neutral (50\%)    & 264 & 75.9\% & 20.3 months & 13 \\
Mild UW (35\%)    &  32 &  9.2\% &  1.6 months &  9 \\
Severe UW (20\%)  &  52 & 14.9\% &  3.5 months & 11 \\
\midrule
\textbf{Total}    & 348 & 100\%  & & \\
\bottomrule
\end{tabular}
\end{table}

The strategy makes on average \textbf{1.62 allocation changes per year},
with the portfolio spending 75.9\% of the time at the neutral 50\%
allocation. The average Neutral dwell of 20.3 months reflects the strategy's
patience: it does not churn in and out at minor signal fluctuations.

\subsection{Performance Summary}

\begin{table}[H]
\centering
\caption{29-year backtest performance (December 1996 -- November 2025)}
\label{tab:perf}
\begin{tabular}{lrrrrr}
\toprule
\textbf{Strategy} & \textbf{Return p.a.} & \textbf{Vol p.a.} & \textbf{Max DD} & \textbf{Sharpe} & \textbf{Avg Eq} \\
\midrule
100\% Equities                & 7.97\% & 15.38\% & $-52.6\%$ & 0.52 & 100.0\% \\
Matched Benchmark (44\% eq)   & 6.04\% &  6.79\% & $-25.1\%$ & 0.89 &  44.1\% \\
\textbf{Two-Step De-Risk + Cash} & \textbf{8.14\%} & \textbf{5.87\%} & $\mathbf{-10.7\%}$ & \textbf{1.39} & \textbf{44.1\%} \\
\bottomrule
\end{tabular}
\end{table}

\noindent
\textbf{Bond rate assumption:} 4.0\% p.a.\ (constant) on the 50\% bond allocation.\\
\textbf{Cash rate:} FRED FEDFUNDS; average over the sample period: \textbf{2.35\% p.a.}

\subsection{Key Observations}

\begin{enumerate}[leftmargin=2em]
  \item \textbf{Return enhancement:} The Two-Step strategy outperforms both
        100\% Equities (+17 bps p.a.) and the Matched Benchmark (+210 bps
        p.a.). The entire excess return over the Matched Benchmark is
        attributable to the timing skill of the MarketComposite signal,
        since both run at identical average equity exposure (44.1\%).

  \item \textbf{Volatility reduction:} Strategy volatility (5.87\%) is
        $62\%$ lower than 100\% Equities (15.38\%) and $14\%$ lower than
        the already de-risked Matched Benchmark (6.79\%). This arises from
        the strategy being underweight equity during high-volatility stress
        episodes.

  \item \textbf{Drawdown control:} The maximum drawdown of $-10.7\%$
        represents a reduction of $80\%$ relative to 100\% Equities
        ($-52.6\%$) and $57\%$ relative to the Matched Benchmark ($-25.1\%$).

  \item \textbf{Sharpe ratio:} The annualised Sharpe of \textbf{1.39}
        compares favourably with 0.89 (Matched Benchmark) and 0.52
        (100\% Equities), demonstrating that the timing signal generates
        positive risk-adjusted alpha.

  \item \textbf{Structural bond allocation:} By holding a constant 50\%
        in bonds, the strategy benefits from bond/equity diversification
        at all times, not just during risk-off episodes. The incremental
        cash position during underweighting additionally reduces
        mark-to-market risk in rising-rate environments.
\end{enumerate}

%% ══════════════════════════════════════════════════════════════════════════
\section{Parameter Reference}
\label{sec:params}
%% ══════════════════════════════════════════════════════════════════════════

\begin{longtable}{p{4.5cm} r p{7.5cm}}
\caption{Complete model parameter reference}\label{tab:params}\\
\toprule
\textbf{Parameter} & \textbf{Default} & \textbf{Description} \\
\midrule
\endfirsthead
\toprule
\textbf{Parameter} & \textbf{Default} & \textbf{Description} \\
\midrule
\endhead
\midrule
\multicolumn{3}{r}{\small(continued on next page)}\\
\endfoot
\bottomrule
\endlastfoot

\multicolumn{3}{l}{\textit{Data parameters}} \\
\texttt{EquityTicker}       & SPY / \texttt{\^{}SPX} & Equity ticker; fallback chain to \texttt{\^{}SPX} \\
\texttt{StartDate}          & 1994-01-01 & Data fetch start (BAMLH0A0HYM2 from 1996; pre-start used for warm-up) \\
\midrule
\multicolumn{3}{l}{\textit{Label thresholds}} \\
\texttt{SevereRetThresh}    & 0 & 3m fwd return below this \% qualifies for Severe \\
\texttt{SevereVolThresh}    & 0.30 & 3m fwd ann vol above this qualifies for Severe \\
\texttt{MildRetThresh}      & 0 & 3m fwd return below this \% qualifies for Mild \\
\midrule
\multicolumn{3}{l}{\textit{Allocation weights}} \\
\texttt{BaseWeight}         & 0.50 & Neutral equity allocation \\
\texttt{MildWeight}         & 0.35 & Mild underweight equity allocation \\
\texttt{SevereWeight}       & 0.20 & Severe underweight equity allocation \\
\midrule
\multicolumn{3}{l}{\textit{EquityTrend signal}} \\
\texttt{MAWindow}           & 200 (days) & Trailing price moving-average window \\
\texttt{TrendScale}         & 10 & tanh scale $k$; 10 $\to$ 10\% above MA $\approx$ 0.76 \\
\midrule
\multicolumn{3}{l}{\textit{PayrollsMom signal}} \\
\texttt{PrefilterWindow}    & 3 (months) & Trailing median noise-gate window \\
\texttt{SmoothHalfLife}     & 3 (months) & EMA half-life for numerator path \\
\texttt{NormWindow}         & 60 (months) & Rolling std window for denominator \\
\texttt{WinsorPct}          & 95 & Expanding winsorisation percentile (denominator) \\
\midrule
\multicolumn{3}{l}{\textit{YieldCurve signal}} \\
\texttt{ZScoreWindow}       & 60 (months) & Trailing z-score window \\
\texttt{SmoothHalfLife}     & 3 (months) & EMA half-life \\
\texttt{TanhScale}          & 1.0 & tanh scale \\
\midrule
\multicolumn{3}{l}{\textit{CreditSpread signal}} \\
\texttt{ZScoreWindow}       & 60 (months) & Trailing z-score window \\
\texttt{SmoothHalfLife}     & 1 (month) & EMA half-life (fast reaction to blowouts) \\
\texttt{TanhScale}          & 1.0 & tanh scale (applied inverted) \\
\texttt{WinsorPct}          & 97 & Expanding winsorisation percentile (denominator) \\
\midrule
\multicolumn{3}{l}{\textit{MarketComposite weights}} \\
\texttt{w\_CreditSpread}    & 0.30 & Weight on CreditSpread component \\
\texttt{w\_EquityTrend}     & 0.50 & Weight on EquityTrend component \\
\texttt{w\_YieldCurve}      & 0.20 & Weight on YieldCurve component \\
\midrule
\multicolumn{3}{l}{\textit{Decision tree}} \\
\texttt{TreeMaxSplits}      & 4 & Maximum tree depth ($\to$ 5 leaves) \\
\texttt{CostSevere}         & 5 & Misclassification cost: Severe predicted as Neutral \\
\texttt{CostMild}           & 2 & Misclassification cost: Mild predicted as Neutral \\
\midrule
\multicolumn{3}{l}{\textit{Portfolio evaluation}} \\
\texttt{BondAnnReturn}      & 0.04 & Assumed constant annual bond return (50\% allocation) \\
\textit{Cash rate}          & FEDFUNDS & FRED effective Federal Funds rate (monthly average) \\

\end{longtable}

%% ══════════════════════════════════════════════════════════════════════════
\section*{Implementation Notes}
\label{sec:impl}
\addcontentsline{toc}{section}{Implementation Notes}
%% ══════════════════════════════════════════════════════════════════════════

\begin{enumerate}[leftmargin=2em]
  \item \textbf{No-timezone datetimes:} Monthly grids use timezone-free
        MATLAB \texttt{datetime} to avoid \texttt{synchronize()} timezone
        conflicts between FRED, Stooq, and Yahoo data.

  \item \textbf{Integer month-matching:} Year$\times$100\,+\,month integer
        keys are used for all monthly alignment operations, providing
        unambiguous month matching across sources with different timestamp
        conventions.

  \item \textbf{Cache system:} \texttt{Model.DataManager} maintains an
        SQLite registry (\texttt{cache/cache\_registry.db}) and per-series
        \texttt{.mat} files. Incremental fetches avoid full re-downloads;
        a fresh request with explicit \texttt{startDate} bypasses the cache
        and fetches the full requested range.

  \item \textbf{UIAxes compatibility:} All chart helpers use \texttt{ax}
        as the first argument and call \texttt{hold(ax,'on')} rather than
        \texttt{axes(ax)}, which is required for MATLAB App Designer
        \texttt{UIAxes} objects.

  \item \textbf{Decision tree availability:} The decision tree approach
        requires the Statistics and Machine Learning Toolbox. If the
        toolbox is absent, \texttt{decisionTreeApproach} raises an error
        that the controller catches gracefully, continuing with threshold
        results only.
\end{enumerate}

%% ── Footer ────────────────────────────────────────────────────────────────
\vfill
\noindent\rule{\linewidth}{0.4pt}
\begin{center}
\small\textcolor{cgray}{%
  CIO Quantitative Research \textbar{} Generated \today \textbar{}
  Backtest period: December 1996 -- November 2025 (29 years) \textbar{}
  Source data: FRED / Stooq
}
\end{center}

\end{document}
